\chapter{Quantum Arithmetic Chern--Simons: Bulk Extension and Residual Obstruction}
\label{v4-arith:w10-a:chapter}

\emph{Chapter~\ref{v4-arith:3d-acs-E1bar} states
\Cref{v4-arith:3d-acs:conj:qu-acs}, the profinite
\(GL_{2}(\widehat{\mathbb Z})\) quantum arithmetic Chern--Simons
theory on \(\overline{\mathrm{Spec}\,\mathbb{Z}}\), as three linked
bulk-boundary axioms. Every formal consequence depends on its
resolution: Ch.~\ref{v4-arith:w9-r1:wilson-frobenius}'s continuous
Wilson-Frobenius identity at \(\varprojlim_{N}^{\mathrm{ren}}\),
Ch.~\ref{v4-arith:w9-r2:archimedean-wilson-surface}'s identification
\(F_{\infty}^{\mathrm{WS}}=F_{\infty}\), the partition formula
\(Z_{\mathrm{ACS}^{\mathrm{qu}}}=\xi\) of
\Cref{v4-arith:3d-acs:cor:partition-bar}, and the triangle of
\Cref{v4-arith:w8-a:cor:triangle}. The bulk measure is examined through
five constructions: (i) state-sum renormalisation along the congruence
tower; (ii) Reshetikhin--Turaev theory for a modular trace lift of the
boundary trace class; (iii) Costello--Gwilliam factorization algebras on the
\'etale site of \(\overline{\mathrm{Spec}\,\mathbb{Z}}\); (iv) Lurie
cobordism hypothesis on an \'etale arithmetic cobordism category;
(v) Scholze six-functor on a conjectural arithmetic diamond. Each
finite-level state sum and the open
\(\mathrm{Spec}\,\mathbb{Z}\setminus\{\infty\}\) factorization algebra
is unconditional. The remaining statements require the archimedean
condensed-Hilbert descent, the modular trace lift, or the arithmetic
diamond.}

\section{Setup and Axiomatic Target}
\label{v4-arith:w10-a:setup}

\begin{definition}[arithmetic 3-space]
\label{v4-arith:w10-a:def:3space}
\ClaimStatusDefinitional. The \emph{arithmetic 3-space} is
\(M=\overline{\mathrm{Spec}\,\mathbb Z}\), the Arakelov
compactification of \(\mathrm{Spec}\,\mathbb Z\) by the archimedean
place \(\{\infty\}\), equipped with its \'etale site
\((M)_{\mathrm{\acute et}}\), its Artin--Verdier fundamental class
\([M]\in H^{3}_{\mathrm{\acute et}}(M,\widehat{\mathbb Z})\), and the
Faltings K\"ahler form \(\omega_{\infty}\) on
\(\Sigma_{\infty}=\{\infty\}\times D^{2}\) of
\Cref{v4-arith:w9-r2:def:Sigma-infty}.
\end{definition}

\begin{definition}[congruence tower and gauge tower]
\label{v4-arith:w10-a:def:tower}
\ClaimStatusDefinitional. The \emph{gauge tower} is the inverse
system of finite groups
\[
\mathcal{T}:=\Bigl\{
G_{N}:=GL_{2}(\mathbb Z/N\mathbb Z),\qquad N\in\mathbb Z_{\geq 1}
\Bigr\}
\]
ordered by divisibility \(M\mid N\) giving reductions \(G_{N}\twoheadrightarrow G_{M}\), with
\[
\varprojlim_{N\in\mathcal{T}}G_{N}\;=\;GL_{2}(\widehat{\mathbb{Z}}).
\]
A \emph{level datum} is a compatible system
\(\alpha_{\bullet}=\{\alpha_{N}\in H^{3}(BG_{N},\mathbb C^{\times})\}_{N}\)
whose inverse-limit class is the arithmetic level
\(k^{\mathrm{Ar}}=\kappa^{\mathrm{Ar}}_{\mathsf G}\) of
Ch.~\ref{v4-arith:3d-acs-E1bar}.
\end{definition}

\begin{remark}[the three axioms recalled]
\label{v4-arith:w10-a:rem:axioms}
\Cref{v4-arith:3d-acs:conj:qu-acs} asks for a renormalised profinite
limit
\[
\mathrm{ACS}^{\mathrm{qu}}_{k^{\mathrm{Ar}}}
:=\varprojlim_{N}^{\mathrm{ren}}Z^{\mathrm{Ar}}_{G_{N},\alpha_{N}}
\]
satisfying:
\textbf{(WZW)}
\(\operatorname{HH}_{0}
(\partial\,\mathrm{ACS}^{\mathrm{qu}}_{k^{\mathrm{Ar}}})
\simeq\mathcal{V}^{\mathrm{prim}}\) on the archimedean boundary, where
\(\mathcal{V}^{\mathrm{prim}}\) is the Hochschild trace class of the
identity on the global arithmetic Iwahori--Hecke category of \(GL_{2}\)
(Ch.~\ref{v4-arith:3d-acs-E1bar});
\textbf{(Z)} \(Z_{\mathrm{ACS}^{\mathrm{qu}}}(\overline{\mathrm{Spec}\,\mathbb Z};k^{\mathrm{Ar}})=\xi(k^{\mathrm{Ar}})\),
where \(\xi\) is the completed Riemann zeta function
(\Cref{v4-arith:3d-acs:cor:partition-bar});
\textbf{(W)} \(\langle W_{\rho}(K_{p})\rangle=\mathrm{tr}\,\rho(\mathrm{Frob}_{p})\),
where \(W_{\rho}(K_{p})\) is the Wilson-loop observable at the prime knot \(K_{p}\subset\overline{\mathrm{Spec}\,\mathbb{Z}}\)
in representation \(\rho\) (Ch.~\ref{v4-arith:w9-r1:wilson-frobenius}).
\end{remark}

\section{Construction I: Tower Renormalisation via Kan Extension}
\label{v4-arith:w10-a:construction-I}

The most direct construction is the one inherited from the finite-level
state-sum construction of
\Cref{v4-arith:3d-acs:thm:finite-level-ACS}. Each
\(Z^{\mathrm{Ar}}_{G_{N},\alpha_{N}}\) is an unconditional finite
state-sum assignment. The obstruction is the existence of a profinite
functor compatible with (WZW), (Z), (W).

\subsection{Finite-level Kan tower}
\label{v4-arith:w10-a:kan-tower}

\begin{lemma}[reduction compatibility at finite levels]
\label{v4-arith:w10-a:lem:reduction}
\ClaimStatusProvedHereConditional. For \(M\mid N\), suppose the kernel
of the reduction map is central abelian and the induced action on the
lifting gerbe is trivial. Then the reduction
\(\pi_{N,M}:G_{N}\twoheadrightarrow G_{M}\) induces a pullback of
level classes
\(\pi_{N,M}^{*}:H^{3}(BG_{M},\mathbb C^{\times})\to H^{3}(BG_{N},\mathbb C^{\times})\)
and a pushforward of field groupoids
\((\pi_{N,M})_{*}:\mathcal F_{G_{N}}(M')\to\mathcal F_{G_{M}}(M')\)
for every arithmetic 3-space \(M'\). With \(\alpha_{N}=\pi_{N,M}^{*}\alpha_{M}\)
consistent, one has
\begin{equation}
\label{v4-arith:w10-a:eq:reduction-compat}
Z^{\mathrm{Ar}}_{G_{N},\alpha_{N}}(M')
\;=\;\bigl|\ker\pi_{N,M}\bigr|^{b_{1}(M')}\cdot Z^{\mathrm{Ar}}_{G_{M},\alpha_{M}}(M')
\end{equation}
where \(b_{1}(M')=\dim_{\mathbb F_{\ell}}H^{1}_{\mathrm{\acute et}}(M',\mathbb F_{\ell})\)
is the first \'etale Betti number with \(\mathbb{F}_{\ell}\)-coefficients for an
auxiliary torsion prime \(\ell\); for \(M'=\overline{\mathrm{Spec}\,\mathbb{Z}}\)
this vanishes by the Minkowski bound, independently of \(\ell\).
\end{lemma}

\begin{proof}
The pullback of cohomology is functorial; the pushforward of field
groupoids is the restriction of scalars
\(\mathcal F_{G_{N}}(M')=\mathcal F_{G_{M}}(M')\times_{\pi_{N,M}}BG_{N}\).
Under the central-abelian hypothesis, the finite-groupoid cardinality
of the kernel of \(\pi_{N,M}\) acting free-transitively on the
\(H^{1}\)-twist gives the displayed factor
\(|\ker\pi_{N,M}|^{b_{1}(M')}\) (cf.\ Freed--Quinn 1993
\emph{Commun.~Math.~Phys.}~156, \S4 for the Dijkgraaf--Witten
version). The level class \(\alpha_{N}=\pi_{N,M}^{*}\alpha_{M}\)
multiplies each field by its \(G_{M}\)-image phase, unchanged by the
torsor of reductions.
\end{proof}

\begin{remark}[divergence is pre-determined]
\label{v4-arith:w10-a:rem:divergence-factor}
Equation \eqref{v4-arith:w10-a:eq:reduction-compat} exposes the
precise source of divergence: the state sum grows by the
volume-factor \(|\ker\pi_{N,M}|^{b_{1}(M')}\) as \(N\uparrow\), so a
naive inverse limit
\(\varprojlim_{N}Z^{\mathrm{Ar}}_{G_{N},\alpha_{N}}\) is ill-defined.
The renormalisation is this volume-factor quotient.
\end{remark}

\begin{definition}[renormalised tower]
\label{v4-arith:w10-a:def:renorm-tower}
\ClaimStatusDefinitional. The \emph{renormalised state-sum tower}
is
\begin{equation}
\label{v4-arith:w10-a:eq:renorm-tower}
\widetilde Z^{\mathrm{Ar}}_{G_{N},\alpha_{N}}(M')
\;:=\;
\bigl|G_{N}\bigr|^{-\chi(M')}\cdot
Z^{\mathrm{Ar}}_{G_{N},\alpha_{N}}(M'),
\end{equation}
where \(\chi(M')\) is the \'etale Euler characteristic of \(M'\) with
\(\mathbb F_{\ell}\)-coefficients, so that the divergence factor
\eqref{v4-arith:w10-a:eq:reduction-compat} is absorbed into the
normalization \(|G_{N}|^{-\chi}\).
\end{definition}

\begin{proposition}[renormalised tower is compatible]
\label{v4-arith:w10-a:prop:renorm-compat}
\ClaimStatusProvedHereConditional. With the renormalisation of
\Cref{v4-arith:w10-a:def:renorm-tower}, for every pair \(M\mid N\)
whose reduction kernel satisfies the central-abelian lifting hypothesis
of \Cref{v4-arith:w10-a:lem:reduction}, the reduction map satisfies
\begin{equation}
\label{v4-arith:w10-a:eq:renorm-compat}
\widetilde Z^{\mathrm{Ar}}_{G_{N},\alpha_{N}}(M')
\;=\;
\widetilde Z^{\mathrm{Ar}}_{G_{M},\alpha_{M}}(M')
\cdot \bigl[G_{N}:G_{M}\bigr]^{-\chi(M')+b_{1}(M')}.
\end{equation}
For \(M'=\overline{\mathrm{Spec}\,\mathbb Z}\), where \(\chi(M')=0\)
by Artin--Verdier duality and \(b_{1}(M')=0\) by the Minkowski bound
on the arithmetic curve, the exponent is zero: the renormalised tower
is a cocycle on such a tower. For the full
\(GL_{2}(\mathbb Z/N)\) congruence tower the nonabelian reduction
kernels require a separate lifting-measure theorem.
\end{proposition}

\begin{proof}
Substitute \eqref{v4-arith:w10-a:eq:reduction-compat} and the
identity
\(|G_{N}|=|G_{M}|\cdot|\ker\pi_{N,M}|=|G_{M}|\cdot[G_{N}:G_{M}]\)
into \eqref{v4-arith:w10-a:eq:renorm-tower}. The ratio
\(\widetilde Z_{N}/\widetilde Z_{M}\) is
\(|G_{N}|^{-\chi}\cdot[G_{N}:G_{M}]^{b_{1}}\cdot|G_{M}|^{\chi}
=[G_{N}:G_{M}]^{-\chi+b_{1}}\),
which is the exponent in \eqref{v4-arith:w10-a:eq:renorm-compat}.
For \(M'=\overline{\mathrm{Spec}\,\mathbb Z}\), \(\chi=0\) by
Artin--Verdier and \(b_{1}=0\) by the cohomological
3-dimensionality (Mazur 1973, Artin--Verdier); the net exponent is
zero.
\end{proof}

\begin{theorem}[renormalised tower admits a limit as a pro-object]
\label{v4-arith:w10-a:thm:profinite-exists}
\ClaimStatusProvedHereConditional. On a reduction tower satisfying the
central-abelian lifting hypothesis at every transition, the renormalised tower
\(\{\widetilde Z^{\mathrm{Ar}}_{G_{N},\alpha_{N}}\}_{N}\) admits a
canonical renormalised profinite limit in the \(\infty\)-category of
pro-objects in the category of state-sum functors on finite
arithmetic cobordism data, given by the \(\infty\)-categorical Kan
right extension along the reduction tower.
\end{theorem}

\begin{proof}
By \Cref{v4-arith:w10-a:prop:renorm-compat}, the renormalised tower
is a levelwise-coherent pro-functor; the formal
\(\infty\)-categorical right Kan extension along the divisibility
tower produces the limit as a pro-object. The finite groupoid
cardinality is preserved under the stated central-abelian transitions;
therefore the corresponding limit is a well-defined element of the
pro-category. The full \(GL_{2}(\widehat{\mathbb{Z}})\) tower is not
included without nonabelian obstruction control.
\end{proof}

\begin{remark}[what Construction I has and has not produced]
\label{v4-arith:w10-a:rem:construction-I-domain}
\Cref{v4-arith:w10-a:thm:profinite-exists} produces the
renormalised inverse limit as a pro-object and confirms the
tower is coherent. This is less than what
\Cref{v4-arith:3d-acs:conj:qu-acs} requires: namely, that the
pro-object \emph{descend to a concrete topological functor} on a
genuine profinite arithmetic cobordism category, and that its
archimedean-boundary Hochschild trace equal \(\mathcal V^{\mathrm{prim}}\).
The descent step is the profinite-measure step. It is not
implied by Kan extension alone.
\end{remark}

\subsection{Sources and Hypotheses for Construction I}
\label{v4-arith:w10-a:ah:construction-I}

The finite-level state sum uses Kim arithmetic Chern--Simons and the
Chern--Dan--Kim--Park non-abelian extension. The renormalisation is the
Dijkgraaf--Witten/Freed--Quinn factor
\(|G_N|^{-\chi}\), and the inverse-limit passage is a Kan extension.
Since \(\chi(\overline{\mathrm{Spec}\,\mathbb Z})=0\), the closed-base
renormalisation factor is \(1\). The limit is a pro-vector-space
object. A topological state functor requires a further condensed or
Tate topology on this pro-object.

\begin{theorem}[Construction I conditional closure]
\label{v4-arith:w10-a:thm:construction-I-closure}
\ClaimStatusProvedHereConditional \emph{(on a choice of condensed
topology on the Kan-limit and its descent to a concrete
topological state-functor)}. Under the profinite-Hilbert descent
hypothesis, the Kan limit of
\Cref{v4-arith:w10-a:thm:profinite-exists} defines
\(\mathrm{ACS}^{\mathrm{qu}}_{k^{\mathrm{Ar}}}\) satisfying the
reduction-coherence portion of
\Cref{v4-arith:3d-acs:conj:qu-acs}. The three axioms (WZW), (Z),
(W) reduce to: (WZW) restriction of the Kan-limit boundary functor
at \(\Sigma_{\infty}\) identifies with
\(\mathcal V^{\mathrm{prim}}\); (Z) evaluation of the Kan-limit on
the closed class \([M]\) reproduces \(\xi(k^{\mathrm{Ar}})\);
(W) evaluation on prime-knot insertions reproduces Frobenius traces.
Each of (WZW), (Z), (W) becomes a \emph{comparison statement} at the
Kan-limit, not a new existence statement.
\end{theorem}

\begin{proof}
The Kan-limit exists by \Cref{v4-arith:w10-a:thm:profinite-exists}.
Under a chosen condensed topology descending the pro-vector-space
structure to a Hilbert-space functor, (WZW), (Z), (W) become
comparison claims at the limit. Each is in principle checkable by
comparing the limit with the right-hand sides of the three axioms
using the finite-level comparisons of Chs.~\ref{v4-arith:w9-r1:wilson-frobenius}
and \ref{v4-arith:3d-acs-E1bar}. The existence of the descent is
hypothetical.
\end{proof}

\section{Construction II: Reshetikhin--Turaev from an Arithmetic Boundary Category}
\label{v4-arith:w10-a:construction-II}

The Reshetikhin--Turaev construction (Turaev 1994
\emph{Quantum Invariants of Knots and 3-Manifolds}, de Gruyter;
Reshetikhin--Turaev 1991 \emph{Invent.~Math.}~103, 547--597)
builds a 3D TQFT from a modular tensor category (MTC) by
state-sum over triangulations. In the arithmetic problem the input
cannot be the category of modules over \(\mathcal V^{\mathrm{prim}}\),
because \(\mathcal V^{\mathrm{prim}}\) is a Hochschild trace class, not
a vertex algebra. The possible Reshetikhin--Turaev input is instead a
modular tensor category whose Hochschild trace and regularised
character have \(\mathcal V^{\mathrm{prim}}\) as their boundary shadow.

\subsection{The arithmetic MTC candidate}
\label{v4-arith:w10-a:arith-MTC}

\begin{definition}[arithmetic modular tensor category candidate]
\label{v4-arith:w10-a:def:arith-MTC}
\ClaimStatusDefinitional. The \emph{arithmetic modular tensor
category candidate} is
\[
\mathcal{MTC}^{\mathrm{Ar}}_{k^{\mathrm{Ar}}}
\;:=\;
\mathcal B^{\mathrm{Ar}}_{k^{\mathrm{Ar}}},
\]
a finite semisimple ribbon category equipped with a trace comparison
\[
        \operatorname{HH}_{0}
        \bigl(\mathcal B^{\mathrm{Ar}}_{k^{\mathrm{Ar}}}\bigr)
        \simeq
        \mathcal V^{\mathrm{prim}}
\]
and with Frobenius-labelled simple objects at the finite places. This
definition is a lifting problem from a Hochschild trace class to a
boundary category; it is not a module category over
\(\mathcal V^{\mathrm{prim}}\).
\end{definition}

\begin{remark}[the modular trace lift is conditional]
\label{v4-arith:w10-a:rem:modularity-cond}
Modularity of \(\mathcal{MTC}^{\mathrm{Ar}}_{k^{\mathrm{Ar}}}\)
requires a non-degenerate braiding, a spherical trace, and compatibility
of the categorical trace with the Iwahori--Hecke Hochschild trace.
The completed-zeta character and finite-prime Fock characters are trace
statements. They do not construct a braided category whose trace is
\(\mathcal V^{\mathrm{prim}}\).
\end{remark}

\subsection{Reshetikhin--Turaev state-sum}
\label{v4-arith:w10-a:RT-state-sum}

\begin{theorem}[Reshetikhin--Turaev from MTC candidate]
\label{v4-arith:w10-a:thm:RT-construction}
\ClaimStatusProvedHereConditional \emph{(on modularity of
\(\mathcal{MTC}^{\mathrm{Ar}}_{k^{\mathrm{Ar}}}\))}. If
\(\mathcal{MTC}^{\mathrm{Ar}}_{k^{\mathrm{Ar}}}\) is modular, then
there is a 3D TQFT functor
\[
Z_{\mathrm{RT}}^{\mathrm{Ar}}:\mathrm{Bord}_{3}^{\mathrm{Ar}}\to\mathrm{Vect}_{\mathbb C},
\]
the Reshetikhin--Turaev state sum on an arithmetic bordism
category \(\mathrm{Bord}_{3}^{\mathrm{Ar}}\) whose objects are
archimedean 2-surfaces and whose morphisms are arithmetic
3-spaces. This functor assigns:
\begin{itemize}
\item a boundary state space whose Hochschild trace is
\(\mathcal V^{\mathrm{prim}}\), after the trace comparison in
\Cref{v4-arith:w10-a:def:arith-MTC};
\item the RT invariant to closed arithmetic 3-spaces;
\item the Wilson-line observable at \(K_{p}\) to the trace of the
corresponding object label on the Frobenius conjugacy class.
\end{itemize}
\end{theorem}

\begin{proof}
Standard Reshetikhin--Turaev construction (Turaev 1994, Ch.~IV):
a modular tensor category produces a 3D TQFT via triangulation.
The boundary state space is computed as morphism spaces in the
MTC; taking the categorical Hochschild trace gives the displayed
boundary trace class. The closed-3-space invariant is
the triangulation-weighted sum over admissible labelings; this
computes the partition function. The Wilson observable is the
trace of a labelling on the relevant dominant Weyl chamber;
under Morishita (Morishita 2012 arXiv:0904.3399) the
dominant-class-function lift of \(\rho(\mathrm{Frob}_{p})\) is the
Frobenius trace.
\end{proof}

\begin{remark}[closure of (WZW) and (W) modulo modularity]
\label{v4-arith:w10-a:rem:RT-closes-WZW-W}
Under modularity of
\(\mathcal{MTC}^{\mathrm{Ar}}_{k^{\mathrm{Ar}}}\) and the trace
comparison of \Cref{v4-arith:w10-a:def:arith-MTC},
\Cref{v4-arith:w10-a:thm:RT-construction} gives the categorical
form of (WZW). The Wilson-line/Frobenius identity is the
Reshetikhin--Turaev quantum-trace formula only after the finite-place
label comparison is supplied. The remaining axiom (Z) is not
automatic: the partition function of an RT TQFT is the
\emph{Reshetikhin--Turaev invariant}, which is a finite complex
number, not an entire function of the level. Matching it with
\(\xi(k^{\mathrm{Ar}})\) requires a separate analytic input.
\end{remark}

\subsection{Partition Identification via the Regularised Boundary Character}
\label{v4-arith:w10-a:zhu-mod}

\begin{proposition}[RT invariant and the regularised boundary character]
\label{v4-arith:w10-a:prop:RT-Zhu}
\ClaimStatusConditional \emph{(on modularity of
\(\mathcal{MTC}^{\mathrm{Ar}}_{k^{\mathrm{Ar}}}\) and on
\Cref{v4-arith:3d-acs:cor:partition-bar}'s regularised-character
map)}. Under modularity, the Reshetikhin--Turaev invariant of
\(\overline{\mathrm{Spec}\,\mathbb{Z}}\) at level \(k^{\mathrm{Ar}}\)
is the regularised trace character of the boundary category. If this
character is identified with the Hochschild trace character of
\(\mathcal{V}^{\mathrm{prim}}\) as in
\Cref{v4-arith:3d-acs:cor:partition-bar}, it is
\(\xi(k^{\mathrm{Ar}})\).
\end{proposition}

\begin{proof}
The Reshetikhin--Turaev invariant of a closed 3-manifold admitting
a Heegaard splitting is the trace of the MTC's canonical endomorphism
(Turaev 1994, Thm.~II.2.3.2). For
\(\overline{\mathrm{Spec}\,\mathbb{Z}}\) viewed as a Heegaard splitting
along \(\Sigma_{\infty}\), this trace is the regularised character of
the boundary category. The comparison with the Hochschild trace class
\(\mathcal V^{\mathrm{prim}}\), and then with \(\xi\), is exactly the
regularised-character input.
\end{proof}

\subsection{Sources and Hypotheses for Construction II}
\label{v4-arith:w10-a:ah:construction-II}

Reshetikhin--Turaev theory supplies a 3D TQFT from a modular tensor
category. Huang--Lepowsky theory is the complex-analytic model in
which a VOA produces such a category. In the arithmetic case the input
is not a module category over \(\mathcal V^{\mathrm{prim}}\), but a
modular trace lift
\(\operatorname{HH}_{0}(\mathcal{MTC}^{\mathrm{Ar}}_{k^{\mathrm{Ar}}})
\simeq\mathcal V^{\mathrm{prim}}\), together with non-degenerate
braiding, spherical trace, finite-place Frobenius labels, and the
regularised character comparison.

\begin{theorem}[Construction II conditional implication]
\label{v4-arith:w10-a:thm:construction-II-closure}
\ClaimStatusProvedHereConditional \emph{(on a modular trace lift,
finite-place Frobenius labels, and the regularised-character
comparison)}. Under these hypotheses, the Reshetikhin--Turaev functor
of \Cref{v4-arith:w10-a:thm:RT-construction} gives the boundary trace
condition, the Wilson-line trace condition, and the scalar partition
identity of \Cref{v4-arith:3d-acs:conj:qu-acs}. Modularity alone does
not imply the zeta partition identity.
\end{theorem}

\begin{proof}
(WZW) is the MTC-boundary identification of
\Cref{v4-arith:w10-a:def:arith-MTC}. (W) is the RT quantum-trace
formula evaluated at the Frobenius conjugacy class via Morishita.
(Z) is \Cref{v4-arith:w10-a:prop:RT-Zhu}. The hypotheses are exactly
the modular trace lift, the finite-place label comparison, and the
regularised-character comparison.
\end{proof}

\section{Construction III: Costello--Gwilliam Factorization on \texorpdfstring{\((\mathrm{Spec}\,\mathbb{Z})_{\mathrm{\acute et}}\)}{the etale site}}
\label{v4-arith:w10-a:construction-III}

The Costello--Gwilliam formalism
(Costello--Gwilliam 2017 \emph{Factorization Algebras in Quantum
Field Theory I}, Cambridge Univ.~Press; vol.~II 2021) builds
perturbative QFTs as factorization algebras on a site. The
\'etale site of \(\mathrm{Spec}\,\mathbb Z\) is a natural substitute
for the usual differential-geometric site in the arithmetic
setting.

\subsection{\'Etale factorization algebra}
\label{v4-arith:w10-a:CG-etale}

\begin{definition}[\'etale factorization prefactor]
\label{v4-arith:w10-a:def:CG-prefactor}
\ClaimStatusDefinitional. Let \(X:=\mathrm{Spec}\,\mathbb{Z}\) (the
open part, without \(\{\infty\}\)). An \emph{\'etale factorization
prefactor} is a prefactorization algebra on the \'etale site
\(X_{\mathrm{\acute et}}\) valued in chain complexes of
\(\mathbb C\)-vector spaces: an assignment
\(U\mapsto\mathcal F(U)\in\mathrm{Ch}_{\mathbb C}\) for each open
\(U\subset X\) in the \'etale topology, together with
quasi-isomorphisms
\(\mathcal F(U_{1})\otimes\cdots\otimes\mathcal F(U_{k})\to\mathcal F(\bigsqcup U_{i})\)
for disjoint unions, compatible with \'etale descent.
\end{definition}

\begin{definition}[finite-level CS prefactor]
\label{v4-arith:w10-a:def:CG-Kim-prefactor}
\ClaimStatusDefinitional. For each finite level \(G_{N}\), the
\emph{Kim prefactor at level \(N\)} is the \'etale prefactorization
algebra
\[
\mathcal F^{\mathrm{Kim}}_{G_{N},\alpha_{N}}(U)
\;:=\;
\mathbb C\bigl[\mathcal F_{G_{N}}(U)\bigr]
\;=\;
\text{the vector space of functions on the finite groupoid}\ \mathcal F_{G_{N}}(U),
\]
with tensor-product structure on disjoint unions given by
restriction of local systems, and
\(\alpha_{N}\)-weighted gluings given by the Kim action
\eqref{v4-arith:3d-acs:eq:kim-action}.
\end{definition}

\begin{theorem}[Kim prefactor is a factorization algebra on open \(X\)]
\label{v4-arith:w10-a:thm:CG-open}
\ClaimStatusProvedHere. The Kim prefactor
\(\mathcal F^{\mathrm{Kim}}_{G_{N},\alpha_{N}}\) of
\Cref{v4-arith:w10-a:def:CG-Kim-prefactor} is a factorization
algebra on the open \'etale site
\((X)_{\mathrm{\acute et}}=(\mathrm{Spec}\,\mathbb{Z})_{\mathrm{\acute et}}\).
Its global sections are
\[
\mathcal F^{\mathrm{Kim}}_{G_{N},\alpha_{N}}(X)
\;=\;
Z^{\mathrm{Ar}}_{G_{N},\alpha_{N}}(X)
\cdot\mathbb C,
\]
the finite-level arithmetic Chern--Simons partition function as
a one-dimensional complex.
\end{theorem}

\begin{proof}
\'Etale descent for local systems with finite fibre (Milne 1980
\emph{Etale Cohomology}, Princeton Univ.~Press, \S II.2) gives the
factorization-algebra structure on \(\mathcal F^{\mathrm{Kim}}\):
for disjoint opens \(U\sqcup V\subset X\),
\(\mathcal F_{G_{N}}(U\sqcup V)=\mathcal F_{G_{N}}(U)\times\mathcal F_{G_{N}}(V)\),
so the algebra of functions factors. The level class \(\alpha_{N}\in H^{3}(BG_{N},\mathbb{Q}/\mathbb{Z})\)
weights each local system via the Kim action
\eqref{v4-arith:3d-acs:eq:kim-action}: for
\(\rho_{N}\colon\pi_{1}^{\mathrm{\acute{e}t}}(X)\to G_{N}\), the phase
is the Artin--Verdier pairing
\(\mathrm{CS}^{\mathrm{Ar}}_{X,\alpha_{N}}(\rho_{N})
=\langle\rho_{N}^{\ast}\alpha_{N},[X]\rangle\in\mathbb{Q}/\mathbb{Z}\),
with \([X]\) the open-part Artin--Verdier fundamental class restricted
from \(\overline{\mathrm{Spec}\,\mathbb{Z}}\) to
\(X=\mathrm{Spec}\,\mathbb{Z}\setminus\{\infty\}\). Fubini for
finite-groupoid cardinality
(\Cref{v4-arith:3d-acs:thm:finite-level-ACS}) gives the global
section identity by summing
\(e^{2\pi i\,\mathrm{CS}^{\mathrm{Ar}}_{X,\alpha_{N}}(\rho_{N})}\)
over isomorphism classes of \(\rho_{N}\).
\end{proof}

\begin{remark}[Construction III closes the finite state-sum \emph{unconditionally}]
\label{v4-arith:w10-a:rem:construction-III-uncond}
On the open part \(X=\mathrm{Spec}\,\mathbb{Z}\setminus\{\infty\}\),
\Cref{v4-arith:w10-a:thm:CG-open} is unconditional. This exposes a
precise structural fact: the Costello--Gwilliam formalism applied
to the open arithmetic curve produces the finite-level Kim ACS
without renormalisation. The divergence of Construction I is localised at
the archimedean boundary.
\end{remark}

\subsection{Archimedean extension}
\label{v4-arith:w10-a:CG-infty}

\begin{proposition}[archimedean factorization extension obstruction]
\label{v4-arith:w10-a:prop:CG-infty-obstruction}
\ClaimStatusProvedHere. The factorization algebra
\(\mathcal F^{\mathrm{Kim}}_{G_{N},\alpha_{N}}\) on the open
\(X=\mathrm{Spec}\,\mathbb{Z}\) does not extend to a factorization
algebra on \(\overline{\mathrm{Spec}\,\mathbb{Z}}\) without an
archimedean insertion: the Arakelov K\"ahler form \(\omega_{\infty}\) on
\(\Sigma_{\infty}\) requires a continuous datum at \(\{\infty\}\) not
supplied by the finite groupoid \(\mathcal F_{G_{N}}(\{\infty\})\).
\end{proposition}

\begin{proof}
The finite field groupoid on the archimedean point
\(\{\infty\}\) is
\(\mathcal F_{G_{N}}(\{\infty\})\simeq BG_{N}\), a discrete classifying
space. The Arakelov K\"ahler form on \(\Sigma_{\infty}\) is a smooth
\((1,1)\)-form on \(D^{2}\) which evaluates over continuous local
systems with K\"ahler weighting; a discrete groupoid
\(BG_{N}\) carries no such structure. The factorization-algebra
extension requires an infinite-dimensional Hilbert space at
\(\Sigma_{\infty}\) (the archimedean Fock) rather than a finite
vector space \(\mathbb C[BG_{N}]\).
\end{proof}

\begin{remark}[where Construction III stops]
\label{v4-arith:w10-a:rem:construction-III-stop}
The Costello--Gwilliam construction closes the finite-prime-only
ACS. The archimedean insertion requires promotion of the
per-place finite groupoid at \(\{\infty\}\) to the
infinite-dimensional archimedean Fock of
Ch.~\ref{v4-arith:w9-r2:archimedean-wilson-surface}. The
obstruction is localised: it is the step from finite vector
spaces to the continuous Hilbert structure at \(\Sigma_{\infty}\).
This is the same obstruction as the Construction I descent from pro-vector
spaces to Hilbert spaces, encountered from a different construction path.
\end{remark}

\subsection{Sources and Hypotheses for Construction III}
\label{v4-arith:w10-a:ah:construction-III}

Costello--Gwilliam factorization algebras, finite \'etale cohomology,
and Kim's arithmetic action give the open-site prefactor. The
\(\alpha_N\)-weight is functorial under \'etale gluing. The object
lives in chain complexes over \(\mathbb C\) on the open
\(\mathrm{Spec}\,\mathbb Z\). Extension to the Arakelov compactification
requires the archimedean Fock factor and the condensed Hilbert
topology at \(\{\infty\}\).

\begin{theorem}[Construction III unconditional closure on the open part]
\label{v4-arith:w10-a:thm:construction-III-closure}
\ClaimStatusProvedHere. On the open arithmetic curve
\(X=\mathrm{Spec}\,\mathbb{Z}\setminus\{\infty\}\), the
Costello--Gwilliam factorization algebra
\(\mathcal F^{\mathrm{Kim}}_{G_{N},\alpha_{N}}\) realizes the
finite-level Kim arithmetic Chern--Simons theory, including the
Wilson-line observable at each prime knot \(K_{p}\subset X\),
unconditionally for every finite \(N\) and every finite level
\(\alpha_{N}\). The profinite limit
\(\varprojlim_{N}^{\mathrm{ren}}\mathcal F^{\mathrm{Kim}}_{G_{N},\alpha_{N}}\)
exists as a pro-object in factorization algebras on \(X_{\mathrm{\acute et}}\)
by \Cref{v4-arith:w10-a:thm:profinite-exists}, and its global
sections compute the profinite partition function on the open
part. The residual obstruction is concentrated at the archimedean
boundary (\Cref{v4-arith:w10-a:prop:CG-infty-obstruction}).
\end{theorem}

\begin{proof}
Each \(\mathcal F^{\mathrm{Kim}}_{G_{N},\alpha_{N}}\) is an
unconditional factorization algebra on \(X_{\mathrm{\acute et}}\) by
\Cref{v4-arith:w10-a:thm:CG-open}. The pro-limit as a pro-object
in factorization algebras is formal; its renormalised limit is
governed by the reduction-compatibility of
\Cref{v4-arith:w10-a:prop:renorm-compat}. The archimedean
extension fails; the obstruction is
\Cref{v4-arith:w10-a:prop:CG-infty-obstruction}.
\end{proof}

\section{Construction IV: Cobordism Hypothesis}
\label{v4-arith:w10-a:construction-IV}

Lurie's cobordism hypothesis (Lurie 2008 arXiv:0905.0465, Thm.~2.4.6)
classifies extended \(n\)-dimensional TQFTs as fully dualizable
objects in a symmetric monoidal \((\infty,n)\)-category. For
\(n=3\) and a target
\(\mathcal{C}=\mathrm{ArithTC}\) the arithmetic tensor category, a
quantum ACS is specified by a fully dualizable object of \(\mathcal{C}\).

\subsection{Fully dualizable object candidate}
\label{v4-arith:w10-a:cobord-candidate}

\begin{definition}[arithmetic target \((\infty,3)\)-category]
\label{v4-arith:w10-a:def:arith-target}
\ClaimStatusDefinitional. The \emph{arithmetic target} is the
symmetric monoidal \((\infty,3)\)-category
\(\mathrm{ArithTC}\) of arithmetic braided tensor categories:
objects are braided tensor categories equipped with a trace functor to
arithmetic Hochschild trace classes; 1-morphisms are bimodule
categories compatible with trace; 2-morphisms are bimodule functors;
3-morphisms are natural transformations; with symmetric monoidal
structure from the Tate restricted tensor of
Ch.~\ref{v4-arith:tate-tensor}.
\end{definition}

\begin{proposition}[fully dualizable object candidate]
\label{v4-arith:w10-a:prop:fully-dualizable-candidate}
\ClaimStatusConditional \emph{(on the finiteness and semisimplicity
of \(\mathcal{MTC}^{\mathrm{Ar}}_{k^{\mathrm{Ar}}}\))}. If
\(\mathcal{MTC}^{\mathrm{Ar}}_{k^{\mathrm{Ar}}}\)
of \Cref{v4-arith:w10-a:def:arith-MTC} is a finite semisimple
braided fusion category with non-degenerate braiding, then it is a
fully dualizable object of \(\mathrm{ArithTC}\).
\end{proposition}

\begin{proof}
Douglas--Schommer-Pries--Snyder 2013 arXiv:1312.7188 classify fully
dualizable objects in the \((\infty,3)\)-category of braided tensor
categories: the fully dualizable objects are exactly the
finite semisimple braided fusion categories with non-degenerate
braiding. Under the modularity hypothesis of
\(\mathcal{MTC}^{\mathrm{Ar}}_{k^{\mathrm{Ar}}}\), the latter is one
such.
\end{proof}

\begin{theorem}[cobordism hypothesis output]
\label{v4-arith:w10-a:thm:cobord-output}
\ClaimStatusProvedHereConditional \emph{(on
\Cref{v4-arith:w10-a:prop:fully-dualizable-candidate}'s modularity
hypothesis)}. If
\(\mathcal{MTC}^{\mathrm{Ar}}_{k^{\mathrm{Ar}}}\) is a modular tensor
category, then the cobordism hypothesis produces an extended 3D TQFT
\[
Z_{\mathrm{coh}}^{\mathrm{Ar}}:\mathrm{Bord}_{3}^{\mathrm{fr}}\to\mathrm{ArithTC}
\]
with
\(Z_{\mathrm{coh}}^{\mathrm{Ar}}(pt)=
\mathcal{MTC}^{\mathrm{Ar}}_{k^{\mathrm{Ar}}}\).
On the value of 2-manifolds and 3-manifolds, this reproduces the
Reshetikhin--Turaev state sum of
\Cref{v4-arith:w10-a:thm:RT-construction}.
\end{theorem}

\begin{proof}
Lurie 2008 (Thm.~2.4.6): the cobordism hypothesis assigns to any
fully dualizable object a unique (up to canonical equivalence)
extended TQFT with that value at a point. Composition with the
evaluation functor on 2- and 3-manifolds reduces the cobordism
assignment to the Reshetikhin--Turaev state sum.
\end{proof}

\begin{remark}[arithmetic cobordism category]
\label{v4-arith:w10-a:rem:arith-cobord}
The cobordism hypothesis applies to a cobordism category
of smooth manifolds. Its arithmetic counterpart requires an
\emph{arithmetic cobordism category} whose objects are
arithmetic 2-surfaces and morphisms are arithmetic 3-spaces.
Construction of such a category is \Cref{v4-arith:3d-acs:residual}
(item 1 of the residual list in Ch.~\ref{v4-arith:3d-acs-E1bar}):
finite arithmetic cobordism functoriality. This is the Construction IV
analogue of Construction III's archimedean extension obstruction.
\end{remark}

\subsection{Sources and Hypotheses for Construction IV}
\label{v4-arith:w10-a:ah:construction-IV}

The cobordism-hypothesis presentation uses Lurie's theorem and the
Douglas--Schommer-Pries--Snyder dualizability criterion for braided
tensor categories. It also uses the Reshetikhin--Turaev state sum on
the resulting modular tensor category. The required inputs are an
arithmetic cobordism category, the modular trace lift, finite-place
label compatibility, and the regularised-character comparison.

\begin{theorem}[Construction IV conditional closure]
\label{v4-arith:w10-a:thm:construction-IV-closure}
\ClaimStatusProvedHereConditional \emph{(on the arithmetic cobordism
category, modular trace lift, finite-place label comparison, and
regularised-character comparison)}. Given these hypotheses, the
cobordism hypothesis produces a unique extended 3D TQFT
\(Z_{\mathrm{coh}}^{\mathrm{Ar}}\) whose trace, Wilson, and scalar
partition data satisfy the axioms of
\Cref{v4-arith:3d-acs:conj:qu-acs}.
\end{theorem}

\begin{proof}
\Cref{v4-arith:w10-a:thm:cobord-output} plus the arithmetic
cobordism extension; the resulting TQFT coincides with the RT
state-sum by uniqueness (Lurie 2008, Cor.~2.4.10).
\end{proof}

\section{Construction V: Six-Functor Formalism on an Arithmetic Diamond}
\label{v4-arith:w10-a:construction-V}

Scholze's six-functor formalism (Scholze 2020--2022
\emph{Lectures on Condensed Mathematics}, Bonn; Scholze 2022
\emph{Six-Functor Formalisms} lecture notes) applies to diamonds and
to a broad class of stacks. A quantum field theory on a space is
given by a ``theory'' in the sense of six-functor formalism:
an association of \(\infty\)-categories with the six operations
\(\otimes,\operatorname{Hom},f^{*},f_{*},f_{!},f^{!}\). For 3D
CS on a diamond, Scholze et al.\ have outlined the six-functor
formalism as the natural home (Fargues--Scholze 2021
arXiv:2102.13459).

\subsection{Diamond of \texorpdfstring{\(\overline{\mathrm{Spec}\,\mathbb Z}\)}{Spec Z bar}}
\label{v4-arith:w10-a:diamond}

\begin{conjecture}[arithmetic diamond]
\label{v4-arith:w10-a:conj:arith-diamond}
\ClaimStatusConjectured. There exists a diamond
\(\overline{\mathrm{Spec}\,\mathbb Z}^{\diamond}\) in the sense of
Scholze 2020 \emph{Etale Cohomology of Diamonds}
(arXiv:1709.07343, \S11) which is a globalisation of the
Fargues--Fontaine curves at each prime, with an archimedean
completion at \(\{\infty\}\) given by a continuous version of
\(\Sigma_{\infty}\). The diamond carries a six-functor formalism
with values in condensed \(\mathbb{Z}_{\ell}\)-sheaves.
\end{conjecture}

\begin{remark}[form of the arithmetic diamond]
\label{v4-arith:w10-a:rem:diamond-form}
\Cref{v4-arith:w10-a:conj:arith-diamond} is a conjecture with no
known construction in the cited literature. The closest object is the
Clausen--Scholze analytic \(\mathbb{Z}\)-spectrum
(Clausen--Scholze 2020 \emph{Condensed Mathematics} lecture notes),
which builds the condensed structure on the arithmetic side but does
not supply the diamond-plus-archimedean-completion we would need.
\end{remark}

\subsection{Six-functor ACS}
\label{v4-arith:w10-a:six-fun-ACS}

\begin{proposition}[six-functor ACS]
\label{v4-arith:w10-a:prop:six-fun-ACS}
\ClaimStatusConditional \emph{(on
\Cref{v4-arith:w10-a:conj:arith-diamond} and on a boundary object
\(\mathscr B_{k^{\mathrm{Ar}}}\) with
\(\operatorname{HH}_{0}(\mathscr B_{k^{\mathrm{Ar}}})
\simeq\mathcal V^{\mathrm{prim}}\))}. If the arithmetic diamond
\(\overline{\mathrm{Spec}\,\mathbb Z}^{\diamond}\) exists with a
six-functor formalism, then the six-functor ACS candidate is
\[
        \mathcal T_{k^{\mathrm{Ar}}}:=
        f^{!}\mathscr B_{k^{\mathrm{Ar}}},
\]
where
\(f:\overline{\mathrm{Spec}\,\mathbb Z}^{\diamond}\to\{pt\}^{\diamond}\)
is the structure map. Its scalar trace is
\[
        \operatorname{Tr}\bigl(f_{!}f^{!}\mathscr B_{k^{\mathrm{Ar}}}\bigr).
\]
Identifying this trace with \(\xi(k^{\mathrm{Ar}})\) is an additional
local-global trace formula, not a formal consequence of the
Fargues--Scholze local geometrisation.
\end{proposition}

\begin{proof}
The six-functor formalism gives the object \(f^{!}\mathscr B\) and its
proper pushforward \(f_{!}f^{!}\mathscr B\). The trace of its identity
is the scalar attached to the closed arithmetic space. Local
Fargues--Scholze geometrisation identifies local spectral-Hecke
categories; passing from those local factors to the completed
Euler product with the archimedean \(\Gamma\)-factor is precisely the
extra local-global trace formula.
\end{proof}

\begin{remark}[Construction V is the most elegant but most conditional]
\label{v4-arith:w10-a:rem:construction-V-form}
Construction V places \Cref{v4-arith:3d-acs:conj:qu-acs} inside
three inputs: the arithmetic diamond, the boundary trace object, and
the local-global trace formula. The diamond is
\Cref{v4-arith:w10-a:conj:arith-diamond}; the trace object lifts
\(\mathcal V^{\mathrm{prim}}\); the trace formula supplies the completed
Euler product.
\end{remark}

\subsection{Sources and Hypotheses for Construction V}
\label{v4-arith:w10-a:ah:construction-V}

The six-functor presentation uses Scholze diamonds, Clausen--Scholze
condensed sheaves, and the Fargues--Scholze spectral category at finite
places. Its inputs are the arithmetic diamond, a boundary trace object
\(\mathscr B_{k^{\mathrm{Ar}}}\), the local-global trace formula, and
finite-place Frobenius compatibility. Proper pushforward and shriek
pullback then give the trace expression
\(\operatorname{Tr}(f_!f^!\mathscr B_{k^{\mathrm{Ar}}})\).

\begin{theorem}[Construction V conditional closure]
\label{v4-arith:w10-a:thm:construction-V-closure}
\ClaimStatusProvedHereConditional \emph{(on the arithmetic diamond,
boundary trace object, local-global trace formula, and finite-place
Frobenius compatibility)}. Under these hypotheses, the six-functor ACS
\(\mathcal T_{k^{\mathrm{Ar}}}=f^{!}\mathscr B_{k^{\mathrm{Ar}}}\)
satisfies (WZW), (Z), (W) of
\Cref{v4-arith:3d-acs:conj:qu-acs}.
\end{theorem}

\begin{proof}
(WZW) is the trace comparison
\(\operatorname{HH}_{0}(\mathscr B_{k^{\mathrm{Ar}}})
\simeq\mathcal V^{\mathrm{prim}}\). (Z) is the local-global trace
formula for
\(\operatorname{Tr}(f_{!}f^{!}\mathscr B_{k^{\mathrm{Ar}}})\). (W) is
the finite-place spectral-Hecke compatibility producing the Frobenius
trace on local Galois representations.
\end{proof}

\section{Construction Comparison and Isolation of the Residual}
\label{v4-arith:w10-a:comparison}

\begin{theorem}[construction comparison]
\label{v4-arith:w10-a:thm:construction-comparison}
\ClaimStatusProvedHere. The five constructions to
\Cref{v4-arith:3d-acs:conj:qu-acs} decompose as follows.
\begin{itemize}
\item Unconditional at finite level: Construction I closes each finite
\(Z^{\mathrm{Ar}}_{G_{N},\alpha_{N}}\); Construction III closes the open-part
factorization algebra.
\item Conditional on a modular trace lift, finite-place Frobenius
labels, and the regularised-character comparison: Constructions II and
IV give Reshetikhin--Turaev and cobordism-hypothesis realisations of
the same boundary, Wilson, and scalar data.
\item Conditional on the arithmetic diamond, boundary trace object,
local-global trace formula, and finite-place Frobenius compatibility:
Construction V gives the six-functor form of the same axioms on
\(\overline{\mathrm{Spec}\,\mathbb Z}^{\diamond}\).
\end{itemize}
\end{theorem}

\begin{proof}
\Cref{v4-arith:w10-a:thm:profinite-exists},
\Cref{v4-arith:w10-a:thm:CG-open},
\Cref{v4-arith:w10-a:thm:construction-II-closure},
\Cref{v4-arith:w10-a:thm:construction-IV-closure},
\Cref{v4-arith:w10-a:thm:construction-V-closure} respectively.
\end{proof}

\subsection{The Common Archimedean Term}
\label{v4-arith:w10-a:common-residual}

\begin{theorem}[archimedean Hilbert descent as common residual condition]
\label{v4-arith:w10-a:thm:common-residual}
\ClaimStatusProvedHere. The finite-level inverse-limit construction,
the open factorization construction, and the six-functor construction
all leave the following archimedean Hilbert descent condition:
\begin{quote}
There exists a choice of condensed (or nuclear or profinite) topology
on the inverse-limit state space
\(\mathcal H^{(\infty),\mathrm{qu}}:=\varprojlim_{N}^{\mathrm{ren}}\mathcal H^{(\infty)}_{G_{N},\alpha_{N}}\)
of finite-level archimedean boundary Hilbert spaces, under which (a)
\(\mathcal H^{(\infty),\mathrm{qu}}\) is the archimedean Fock factor
\(\mathcal V^{\mathrm{prim}}_{\mathrm{Heis},\infty}\); and (b) the
Arakelov-K\"ahler Wilson-surface observable \(F_{\infty}^{\mathrm{WS}}\)
of \Cref{v4-arith:w9-r2:def:F-infty-WS} acts on this inverse limit
as the scaling Frobenius \(F_{\infty}=e^{-\Theta_{\infty}}
=e^{-\frac12\mathrm{id}-i\widehat D}\) of
\Cref{v4-arith:w9-r2:def:F-infty}.
\end{quote}
Equivalently, the residual asks that the archimedean K\"ahler form
\(\omega_{\infty}\) induce a condensed-limit action of
\(\mathbb{R}^{\times}_{+}\) on the limit state space, and that this
action be unitarily conjugate to the Weil--Petersson scaling
representation on \(\mathcal V^{\mathrm{prim}}_{\mathrm{Heis},\infty}\).
The Reshetikhin--Turaev and cobordism-hypothesis constructions add a
separate modular trace lift of
\(\mathcal V^{\mathrm{prim}}\); this lift is sufficient for the
boundary category but is not identical to the Hilbert descent.
\end{theorem}

\begin{proof}
Construction I reduces to this descent in
\Cref{v4-arith:w10-a:rem:construction-I-domain}. Construction III reduces to it
in \Cref{v4-arith:w10-a:rem:construction-III-stop}. Construction V
requires the arithmetic diamond, whose six-functor formalism would
contain the descent in its \(f^{!}\) on the archimedean completion.
Constructions II and IV instead require a modular tensor category
lifting the same trace class. The lifting condition maps to the
Hilbert descent through the boundary trace but is a stronger
categorical input, not a formal consequence.
\end{proof}

\begin{remark}[the residual has three aliases]
\label{v4-arith:w10-a:rem:aliases}
The archimedean Hilbert descent of
\Cref{v4-arith:w10-a:thm:common-residual} is the same obstruction as:
\begin{itemize}
\item \emph{Construction I alias}: profinite-Hilbert descent of the Kan
inverse limit.
\item \emph{Construction III alias}: extension of the Costello--Gwilliam
factorization algebra from the open \'etale curve to the full
Arakelov compactification.
\item \emph{Construction V alias}: archimedean completion of the arithmetic
diamond with its induced condensed structure.
\end{itemize}
These are three presentations of the fact that the
\'etale-discrete state spaces at finite level must descend to a
continuous Hilbert space at the archimedean boundary under the
Arakelov K\"ahler form.
\end{remark}

\subsection{Relation to the Hilbert--P\'olya/Connes--Deninger obstruction}
\label{v4-arith:w10-a:HP-relation}

\begin{remark}[archimedean descent is strictly weaker than RH]
\label{v4-arith:w10-a:rem:weaker-than-RH}
The archimedean Hilbert descent of
\Cref{v4-arith:w10-a:thm:common-residual} is an \emph{existence}
claim for the state space \(\mathcal H^{(\infty),\mathrm{qu}}\) and a
\emph{functoriality} claim for the Arakelov scaling action. It is
strictly weaker than the Hilbert--P\'olya hypothesis
H-P\(^{\mathrm{Ar}}\) of \Cref{v4-arith:rh:hyp:HPAr}: H-P\(^{\mathrm{Ar}}\)
requires, in addition, that
\(\Theta_{\xi}=\frac12\mathrm{id}+iH_{\mathrm{HP}}\) on
\(H^{1}_{\mathrm{ch}}\subset\mathcal H^{(\infty),\mathrm{qu}}\) have
regularised determinant divisor \(\mathrm{div}\,\xi_{\mathbb R}\) and
Weil--Connes trace distribution. The archimedean Hilbert
descent does not assert this; it only constructs the state space
and the scaling action.
\end{remark}

\begin{theorem}[Clausen--Scholze condensed-Hilbert descent plus A-TP \(\Rightarrow\) RH]
\label{v4-arith:w10-a:thm:descent-plus-ATP-equals-RH}
\ClaimStatusProvedHereConditional \emph{(on the archimedean
condensed-Hilbert descent of
\Cref{v4-arith:w10-a:conj:condensed-descent} and the
A-TP archimedean positivity of
Ch.~\ref{v4-arith:w5-a:conj:archimedean-tate-positivity})}. Under
both the Clausen--Scholze condensed-Hilbert descent and A-TP, the Riemann
Hypothesis holds.
\end{theorem}

\begin{proof}
The condensed-Hilbert descent constructs
\(\mathcal H^{(\infty),\mathrm{qu}}\) on which
\(F_{\infty}=e^{-\Theta_{\infty}}=e^{-1/2-i\widehat D}\) acts, with
the Tate half-shift included. By
\Cref{v4-arith:w9-r2:thm:purity-ATP-RH}, A-TP is
equivalent to \(F_{\infty}\)-weight \(=1/2\), which is equivalent to RH.
The descent supplies the Hilbert realisation of the scaling operator.
The zero-placement assertion enters only through A-TP; it is not a
spectral consequence of the descent.
\end{proof}

\begin{remark}[the archimedean Hilbert descent does not imply RH alone]
\label{v4-arith:w10-a:rem:what-resolves}
Resolving the archimedean Hilbert descent alone does not
resolve RH: it constructs
\(\mathrm{ACS}^{\mathrm{qu}}_{k^{\mathrm{Ar}}}\) only on the
descended Hilbert-trace clauses (WZW), (Z), and (W); it does not
supply the full cobordism functor, nor A-TP positivity
of the scaling operator. Under both descent and A-TP, RH follows.
The archimedean Hilbert descent is the bulk-construction
obstruction; A-TP is the separate positivity obstruction. After the
descent, the remaining bulk gluing assertions stay separate; the
RH-equivalent condition in this construction is A-TP.
\end{remark}

\section{Refinement: The Condensed-Hilbert Bulk}
\label{v4-arith:w10-a:condensed-bulk}

\subsection{Condensed-Hilbert setup}
\label{v4-arith:w10-a:condensed-setup}

\begin{definition}[condensed Hilbert pro-system]
\label{v4-arith:w10-a:def:condensed-pro-system}
\ClaimStatusDefinitional. The \emph{condensed Hilbert pro-system}
at the archimedean boundary is the pro-system
\[
\mathcal H^{(\infty)}_{\bullet}\;:=\;
\bigl\{
\mathcal H^{(\infty)}_{G_{N},\alpha_{N}}
\;=\;\text{finite-level boundary state space from}\ \Cref{v4-arith:3d-acs:def:finite-boundary-category}
\bigr\}_{N},
\]
viewed as a pro-object in the symmetric monoidal
\(\infty\)-category \(\mathrm{Cond}_{\mathrm{Hilb}}(\mathbb C)\) of
condensed Hilbert spaces in the sense of Clausen--Scholze 2020,
with structure maps given by the reduction \(\pi_{N,M}\) and a
chosen renormalisation weight absorbing the \(|G_{N}|^{-\chi}\)
factor of \Cref{v4-arith:w10-a:def:renorm-tower}.
\end{definition}

\begin{definition}[target archimedean Hilbert]
\label{v4-arith:w10-a:def:target-Hilbert}
\ClaimStatusDefinitional. The \emph{target archimedean Hilbert
space} is
\(\mathcal V^{\mathrm{prim}}_{\mathrm{Heis},\infty}\), the archimedean
Fock factor of the arithmetic Heisenberg algebra
(Ch.~\ref{v4-arith:virasoro-bootstrap}), viewed as a condensed
Hilbert space with its Arakelov/Weil--Petersson inner product
(Ch.~\ref{v4-arith:w5-a:chapter}).
\end{definition}

\subsection{The Condensed Hilbert--Trace Package}
\label{v4-arith:w10-a:named-obstruction}

\begin{conjecture}[archimedean condensed-Hilbert descent and trace]
\label{v4-arith:w10-a:conj:condensed-descent}
\ClaimStatusConjectured. The condensed-Hilbert pro-system
\(\mathcal H^{(\infty)}_{\bullet}\) of
\Cref{v4-arith:w10-a:def:condensed-pro-system} admits a limit
\(\mathcal H^{(\infty)}_{\infty}\) in \(\mathrm{Cond}_{\mathrm{Hilb}}(\mathbb C)\),
and there is an isomorphism of condensed Hilbert spaces
\begin{equation}
\label{v4-arith:w10-a:eq:condensed-descent}
\mathcal H^{(\infty)}_{\infty}
\;\cong\;
\mathcal V^{\mathrm{prim}}_{\mathrm{Heis},\infty}
\end{equation}
intertwining the Arakelov scaling \(\mathbb{R}^{\times}_{+}\)-action
on the target with the Wilson-surface
\(F_{\infty}^{\mathrm{WS}}\)-action on the limit.
In addition, the regularised trace of the descended Kan-limit is the
regularised trace character of
\(\mathcal V^{\mathrm{prim}}_{\mathrm{Heis},\infty}\), hence
\(\xi(k^{\mathrm{Ar}})\), and the finite-place Wilson endomorphisms
descend to the Frobenius traces
\(\mathrm{tr}\,\rho(\mathrm{Frob}_{p})\).
\end{conjecture}

\begin{theorem}[condensed Hilbert--trace package supplies the matching ACS clauses]
\label{v4-arith:w10-a:thm:descent-equivalent}
\ClaimStatusProvedHereConditional \emph{(after choosing the same
bulk-boundary ACS model)}. \Cref{v4-arith:w10-a:conj:condensed-descent}
supplies the archimedean Hilbert, trace, and Wilson clauses of
\Cref{v4-arith:3d-acs:conj:qu-acs}. It is not, by itself, the full
topological ACS functor.
\end{theorem}

\begin{proof}
Given \Cref{v4-arith:w10-a:conj:condensed-descent},
the limit \(\mathcal H^{(\infty)}_{\infty}\) is the archimedean
boundary Hilbert space of the profinite ACS; combined with the
Construction III factorization algebra on the open curve, this descends
to a functorial assignment on
\(\overline{\mathrm{Spec}\,\mathbb{Z}}\) satisfying (WZW). The
partition function (Z) is the regularised trace clause. The Wilson-line
axiom (W) is the finite-place descent clause.
The remaining assertion of \Cref{v4-arith:3d-acs:conj:qu-acs} is the
existence of the full cobordism functor with gluing; it is not implied
by the condensed Hilbert limit alone.
\end{proof}

\begin{remark}[form of the condensed statement]
\label{v4-arith:w10-a:rem:named-theorem}
\Cref{v4-arith:w10-a:conj:condensed-descent} is a condensed Hilbert
limit together with a trace identity and finite-place endomorphism
compatibility. The Hilbert limit alone gives only the archimedean
state space; the trace and Wilson clauses carry the scalar and
finite-prime content of the quantum ACS axioms.
\end{remark}

\section{Wilson-Observable Compatibility}
\label{v4-arith:w10-a:wilson-compat}

\begin{theorem}[Wilson-line compatibility]
\label{v4-arith:w10-a:thm:wilson-line-compat}
\ClaimStatusProvedHereConditional \emph{(on
\Cref{v4-arith:w10-a:conj:condensed-descent})}. Under the
condensed-Hilbert descent, the Wilson-line observable at a prime
knot \(K_{p}\) in the bulk \(\mathrm{ACS}^{\mathrm{qu}}_{k^{\mathrm{Ar}}}\),
evaluated on the limit state space
\(\mathcal H^{(\infty)}_{\infty}\cong\mathcal V^{\mathrm{prim}}_{\mathrm{Heis},\infty}\),
is exactly the Frobenius trace \(\mathrm{tr}\,\rho(\mathrm{Frob}_{p})\)
of \Cref{v4-arith:w9-r1:thm:main}(2), as a Wilson clause of the
descent.
\end{theorem}

\begin{proof}
\Cref{v4-arith:w10-a:thm:descent-equivalent} supplies the Wilson
clause of the chosen descended ACS model. Axiom (W) in that model is
the Wilson-line Frobenius identity. The finite-level identity
\Cref{v4-arith:w9-r1:thm:main}(1) is unconditional; the profinite
limit \Cref{v4-arith:w9-r1:thm:main}(2) is conditional on the
renormalised profinite limit, which is supplied by the descent.
\end{proof}

\begin{theorem}[Wilson-surface compatibility]
\label{v4-arith:w10-a:thm:wilson-surface-compat}
\ClaimStatusProvedHereConditional \emph{(on
\Cref{v4-arith:w10-a:conj:condensed-descent})}. Under the
condensed-Hilbert descent, the archimedean Wilson-surface
observable \(F_{\infty}^{\mathrm{WS}}\) of
\Cref{v4-arith:w9-r2:def:F-infty-WS}, acting on the limit state
space \(\mathcal H^{(\infty)}_{\infty}\), coincides with the Arakelov
scaling Frobenius \(F_{\infty}=e^{-\Theta_{\infty}}\) of
\Cref{v4-arith:w9-r2:def:F-infty}, as the archimedean Wilson clause
of the descent.
\end{theorem}

\begin{proof}
Immediate from \Cref{v4-arith:w9-r2:thm:WS-equals-scaling}: the
condensed-Hilbert descent supplies the archimedean Wilson-surface
clause needed there. It does not replace the remaining cobordism
gluing clauses of the full quantum ACS conjecture.
\end{proof}

\begin{corollary}[Wilson observables on the descended model]
\label{v4-arith:w10-a:cor:wilson-one-residual}
\ClaimStatusProvedHereConditional \emph{(on
\Cref{v4-arith:w10-a:conj:condensed-descent})}. The Wilson-line Frobenius identity
(Ch.~\ref{v4-arith:w9-r1:wilson-frobenius}) and the archimedean
Wilson-surface = Arakelov scaling identity
(Ch.~\ref{v4-arith:w9-r2:archimedean-wilson-surface}) are the two
Wilson clauses of the archimedean condensed-Hilbert descent.
\end{corollary}

\begin{proof}
\Cref{v4-arith:w10-a:thm:wilson-line-compat} and
\Cref{v4-arith:w10-a:thm:wilson-surface-compat}.
\end{proof}

\section{Comparison of the Five Presentations}
\label{v4-arith:w10-a:ah:synthesis}

\Cref{v4-arith:w10-a:thm:common-residual} compares five independent
presentations. The finite state-sum presentation uses
Dijkgraaf--Witten renormalisation in the form of Freed--Quinn. The
categorical presentation uses the Reshetikhin--Turaev state sum. The
local-field presentation uses Costello--Gwilliam factorization on the
\'etale site. The fully extended presentation uses the cobordism
hypothesis and the Douglas--Schommer-Pries--Snyder dualizability
criterion. The six-functor presentation uses the conjectural
arithmetic diamond and the Fargues--Scholze spectral side.

The Euler characteristic
\(\chi(\overline{\mathrm{Spec}\,\mathbb Z})=0\) makes the finite
renormalisation factor \(|G_N|^{-\chi}\) trivial on the closed base.
The non-trivial step is not a sign convention; it is the passage
\[
\mathrm{Vect}_{\mathbb C}
\longrightarrow \mathrm{Pro}(\mathrm{Vect}_{\mathbb C})
\longrightarrow \mathrm{Cond}_{\mathrm{Hilb}}(\mathbb C)
\]
with the Arakelov K\"ahler scaling acting unitarily at the limit.
This is precisely the Hilbert--trace descent of
\Cref{v4-arith:w10-a:conj:condensed-descent}. The
Reshetikhin--Turaev and cobordism-hypothesis presentations further
require the modular trace lift; the diamond presentation further
requires the boundary trace object and the local-global trace formula.

\section{Remaining Inputs}
\label{v4-arith:w10-a:residual}

Three inputs remain.

\begin{enumerate}
\item \emph{Archimedean condensed-Hilbert descent}
(\Cref{v4-arith:w10-a:conj:condensed-descent}). This is the precise
Hilbert, trace, and Wilson matching statement compared with
\Cref{v4-arith:3d-acs:conj:qu-acs} by
\Cref{v4-arith:w10-a:thm:descent-equivalent}. Domain: a
condensed-mathematics limit problem in the Clausen--Scholze
framework. Well-defined, unresolved.

\item \emph{Arithmetic modular trace lift}
(\Cref{v4-arith:w10-a:rem:modularity-cond}). This is a sufficient
condition for Constructions II and IV. It is strictly stronger than
Residual 1: it asks for a modular tensor category whose Hochschild
trace is \(\mathcal V^{\mathrm{prim}}\), with finite-place Frobenius
labels and a regularised character. Huang--Lepowsky theory is the
complex-analytic model for how a VOA can produce an MTC; here the
problem is a trace-lifting problem, since
\(\mathcal V^{\mathrm{prim}}\) is already a trace class.

\item \emph{Arithmetic diamond}
(\Cref{v4-arith:w10-a:conj:arith-diamond}). This supplies the
six-functor space for Construction V. It must be joined to a boundary
trace object and to the local-global trace formula; the diamond alone
does not force the zeta partition identity.
\end{enumerate}

\begin{remark}[relations among the three inputs]
\label{v4-arith:w10-a:rem:partial-order}
The modular trace lift and the arithmetic diamond are sufficient
geometrisations of the boundary problem. Neither one is a formal
consequence of the other. Both must still meet the same archimedean
descent and local-global trace constraints before the
\(\mathrm{ACS}^{\mathrm{qu}}\) axioms hold.
\end{remark}

\section{Unconditional Pro-Functorial Core and Residual Boundary}
\label{v4-arith:w10-a:what-closes}

\subsection*{Unconditional consequences}

\begin{enumerate}
\item The finite-level renormalised state-sum tower is a coherent
pro-functor; its Kan limit exists as a pro-object
(\Cref{v4-arith:w10-a:thm:profinite-exists}).
\item The Costello--Gwilliam factorization algebra on the open
\'etale curve \(\mathrm{Spec}\,\mathbb{Z}\setminus\{\infty\}\) is an
unconditional factorization algebra for every finite level
(\Cref{v4-arith:w10-a:thm:CG-open}).
\item The construction comparison theorem
(\Cref{v4-arith:w10-a:thm:construction-comparison}) and the
archimedean-descent theorem
(\Cref{v4-arith:w10-a:thm:common-residual}) isolate the common
Hilbert-space term in the finite-level, factorization, and six-functor
forms.
\item The descent comparison theorem
(\Cref{v4-arith:w10-a:thm:descent-equivalent}) shows that the
condensed-Hilbert descent supplies the Hilbert, trace, and Wilson
clauses of the chosen ACS model; the full cobordism functor and gluing
remain separate.
\item Under the descent, the Wilson-line and Wilson-surface
identities hold inside the descended model
(\Cref{v4-arith:w10-a:cor:wilson-one-residual}).
\end{enumerate}

\subsection*{Conditional on the condensed-Hilbert descent}

\begin{enumerate}
\item The three axioms (WZW), (Z), (W) of
\Cref{v4-arith:3d-acs:conj:qu-acs}.
\item Ch.~\ref{v4-arith:w9-r1:wilson-frobenius} profinite
Wilson-Frobenius identity.
\item Ch.~\ref{v4-arith:w9-r2:archimedean-wilson-surface}
Wilson-surface = Arakelov scaling.
\item The triangle of \Cref{v4-arith:w8-a:cor:triangle} (partition
function = zeta character = regularised character of candidate
Deninger complex).
\end{enumerate}

\subsection*{Still strictly RH-level beyond the descent}

\begin{enumerate}
\item The A-TP archimedean positivity of
Ch.~\ref{v4-arith:w5-a:chapter} (equivalent to RH by
\Cref{v4-arith:w9-r2:thm:purity-ATP-RH}).
\item The Hilbert--P\'olya determinant-trace realisation H-P\(^{\mathrm{Ar}}\)
of \Cref{v4-arith:rh:hyp:HPAr} (equivalent to zero-placement
determinant identity).
\item The Deninger spectral-flow hypothesis
\Cref{v4-arith:3d-acs:conj:spectral-flow} (equivalent to the full
Weil--Guinand distribution).
\end{enumerate}

The chapter does not reduce the RH-level obstructions. It identifies
the condensed-Hilbert descent as the common archimedean term in the
bulk comparison.

\section{Convergence}
\label{v4-arith:w10-a:summary}

The finite-level, factorization, and six-functor forms of
\(\mathrm{ACS}^{\mathrm{qu}}_{k^{\mathrm{Ar}}}\) contain the
archimedean condensed-Hilbert descent of
\Cref{v4-arith:w10-a:conj:condensed-descent}. Constructions I and III close
unconditionally on the finite-level and open-curve parts; Constructions
II and IV require a modular trace lift with the regularised character
comparison, while Construction V requires the arithmetic diamond
together with the boundary trace object and the local-global trace
formula.
Under the descent, both Wilson observables are defined in the
descended model; under the descent together with A-TP, the Riemann
Hypothesis holds
(\Cref{v4-arith:w10-a:thm:descent-plus-ATP-equals-RH}).

The archimedean fibre is the essential locus: finite-level
\'etale-discrete state spaces descend to
discrete finite-prime factorization algebras unconditionally; the
archimedean Hilbert structure is where the descent either succeeds
or fails. This is the same concentration phenomenon
as in Ch.~\ref{v4-arith:w9-r2:archimedean-wilson-surface}, where
the Riemann Hypothesis concentrates in the archimedean fibre. Both
point to the same mathematical joint: the arithmetic K\"ahler form
\(\omega_{\infty}\) on \(\Sigma_{\infty}\) is the carrier of the
analytic content that the \'etale-finite theory cannot supply.

\section{Bibliography}
\label{v4-arith:w10-a:bibliography}

\begin{description}
\item[Clausen--Scholze 2020.] Clausen, D.; Scholze, P. \emph{Condensed
Mathematics}. Lecture notes, Bonn; available at Scholze's homepage.
\item[Costello--Gwilliam 2017.] Costello, K.; Gwilliam, O.
\emph{Factorization Algebras in Quantum Field Theory}, Vol.~I.
Cambridge Univ.~Press, 2017.
\item[Costello--Gwilliam 2021.] Costello, K.; Gwilliam, O.
\emph{Factorization Algebras in Quantum Field Theory}, Vol.~II.
Cambridge Univ.~Press, 2021.
\item[Douglas--Schommer-Pries--Snyder 2013.] Douglas, C.~L.;
Schommer-Pries, C.~J.; Snyder, N. \emph{Dualizable tensor
categories}. arXiv:1312.7188.
\item[Fargues--Scholze 2021.] Fargues, L.; Scholze, P.
\emph{Geometrization of the local Langlands correspondence}.
arXiv:2102.13459.
\item[Freed--Quinn 1993.] Freed, D.~S.; Quinn, F. \emph{Chern--Simons
theory with finite gauge group}. Commun.~Math.~Phys.~\textbf{156}
(1993), 435--472.
\item[Huang--Lepowsky 1994.] Huang, Y.-Z.; Lepowsky, J.
\emph{Tensor products of modules for a vertex operator algebra and
vertex tensor categories}. In: \emph{Lie Theory and Geometry}
(Brylinski et al., eds.), Birkhauser, 1994, 349--383.
\item[Kim 2015.] Kim, M. \emph{Arithmetic Chern--Simons theory I}.
arXiv:1510.05818.
\item[Lurie 2008.] Lurie, J. \emph{On the classification of
topological field theories}. arXiv:0905.0465.
\item[Lurie HTT.] Lurie, J. \emph{Higher Topos Theory}. Princeton
Univ.~Press, 2009.
\item[Milne 1980.] Milne, J.~S. \emph{Etale Cohomology}. Princeton
Univ.~Press, 1980.
\item[Reshetikhin--Turaev 1991.] Reshetikhin, N.; Turaev, V.~G.
\emph{Invariants of 3-manifolds via link polynomials and quantum
groups}. Invent.~Math.~\textbf{103} (1991), 547--597.
\item[Scholze 2017.] Scholze, P. \emph{Etale cohomology of diamonds}.
arXiv:1709.07343.
\item[Scholze 2020.] Scholze, P. \emph{Six-functor formalism}.
Lecture notes, Bonn.
\item[Turaev 1994.] Turaev, V.~G. \emph{Quantum Invariants of Knots
and 3-Manifolds}. de Gruyter, 1994.
\end{description}
